\section{Erd\H{o}s--Hajnal reciprocal cycle-length sum (Problem 65) --- Round 6}

\subsection*{1) ROUND-6 OBJECTIVE}

\textbf{Path chosen: (A) proof-oriented gap closure.}
Round~5 proved that on the complete-bipartite line
\[
|E(G)|=2|V(G)|-4,
\]
one has
\[
\min S(G)=\frac14
\qquad\text{for all }7\le n\le 23,
\]
with equality attained by $K_{2,n-2}$, and identified $n=24$ as the first place where the old ``force $C_5$ and $C_6$'' argument stops.

In this round I keep the same target family but change the forcing pair: instead of trying to force $C_5$ and $C_6$, I force $C_6$ and $C_8$.
This yields a substantial extension:
\[
\min\{S(G): |V(G)|=n,\ |E(G)|=2n-4\}=\frac14
\qquad\text{for every }7\le n\le 39,
\]
again with extremal graph $K_{2,n-2}$.

\subsection*{2) Round-5 FOUNDATION USED}

I rely on the following vetted Round~5 results without re-proving them.

\begin{itemize}
\item[(F1)] For $2\le a\le b$, the complete bipartite graph $K_{a,b}$ has
\[
\mathcal C(K_{a,b})=\{4,6,\dots,2a\},
\qquad
S(K_{a,b})=\sum_{j=2}^a\frac1{2j}=\frac12(H_a-1).
\]
In particular,
\[
S(K_{2,n-2})=\frac14.
\]

\item[(F2)] Round~5 theorem: for every $7\le n\le 23$,
\[
\min\{S(G): |V(G)|=n,\ |E(G)|=2n-4\}=\frac14,
\]
attained by $K_{2,n-2}$.
\end{itemize}

\subsection*{3) NEW INSIGHT / TOOL (ROUND-6)}

The new idea is to replace the Round~5 forcing pair $(C_5,C_6)$ by the pair $(C_6,C_8)$.
The relevant exact extremal numbers are
\[
\alpha_n:=\operatorname{ex}(n;\{C_3,C_4,C_6\}),
\qquad
\beta_n:=\operatorname{ex}(n;\{C_4,C_8\}).
\]
From the exact ANU tables one has, for $24\le n\le 39$,
\[
\begin{array}{c|cccccccccccccccc}
 n & 24&25&26&27&28&29&30&31&32&33&34&35&36&37&38&39\\\hline
 \alpha_n & 38&40&42&45&46&48&51&53&55&57&60&62&65&67&69&72\\
 \beta_n  & 38&40&42&44&46&48&50&52&54&56&58&59&61&63&65&67\\
 2n-4      & 44&46&48&50&52&54&56&58&60&62&64&66&68&70&72&74
\end{array}
\]
so throughout this whole range,
\[
2n-4>\alpha_n
\qquad\text{and}\qquad
2n-4>\beta_n.
\]
Hence any $n$-vertex graph with $2n-4$ edges and no $C_3,C_4$ must contain a $C_6$; and if it also has no $C_5$, then the absence of a $C_4$ plus the inequality $2n-4>\beta_n$ forces a $C_8$.

This is genuinely new beyond Round~5, because $C_5$ is no longer forceable from $n=24$ onward, while $C_8$ still is.

\subsection*{4) ATTACK PLAN (ROUND-6)}

The exact gap left by Round~5 was the unresolved tail of the line $m=2n-4$ beyond $n=23$.

The Round~6 plan is:
\begin{enumerate}
\item Use the exact values of $\alpha_n=\operatorname{ex}(n;\{C_3,C_4,C_6\})$ to force a $6$-cycle whenever $24\le n\le 39$.
\item Use the exact values of $\beta_n=\operatorname{ex}(n;\{C_4,C_8\})$ to force an $8$-cycle whenever $24\le n\le 39$ and $G$ has no $4$-cycle.
\item Split into cases:
  \begin{itemize}
  \item if $G$ has $C_3$ or $C_4$, then automatically $S(G)\ge 1/4$;
  \item if $G$ has no $C_3,C_4$ but has a $C_5$, then together with the forced $C_6$ we get
  \[
  S(G)\ge \frac15+\frac16>\frac14;
  \]
  \item if $G$ has no $C_3,C_4,C_5$, then the forced $C_6$ and $C_8$ give
  \[
  S(G)\ge \frac16+\frac18=\frac7{24}>\frac14.
  \]
  \end{itemize}
\item Compare with $K_{2,n-2}$, which always gives equality $S=1/4$.
\end{enumerate}

\subsection*{5) WORK (ROUND-6)}

As before, for a graph $G$ let
\[
\mathcal C(G):=\{\ell\ge 3 : \text{$G$ contains a cycle of length $\ell$}\},
\qquad
S(G):=\sum_{\ell\in\mathcal C(G)}\frac1\ell.
\]

\begin{theorem}[Extension of the $K_{2,n-2}$ line through $n=39$]\label{thm:round6-main}
For every integer $7\le n\le 39$,
\[
\min\{S(G): |V(G)|=n,\ |E(G)|=2n-4\}=\frac14.
\]
The minimum is attained by the complete bipartite graph $K_{2,n-2}$.
\end{theorem}

\begin{proof}
By (F2), the statement is already known for $7\le n\le 23$.
So it remains only to consider $24\le n\le 39$.

Fix such an $n$, and let $G$ be any graph on $n$ vertices with $2n-4$ edges.
By (F1), the complete bipartite graph $K_{2,n-2}$ has only $4$-cycles, so
\[
S(K_{2,n-2})=\frac14.
\]
Thus it suffices to prove that every such $G$ satisfies $S(G)\ge \frac14$.

We split into cases.

\smallskip
\noindent\textbf{Case 1: $G$ contains a triangle $C_3$.}
Then
\[
S(G)\ge \frac13>\frac14.
\]

\smallskip
\noindent\textbf{Case 2: $G$ contains a $4$-cycle $C_4$.}
Then
\[
S(G)\ge \frac14.
\]

\smallskip
\noindent\textbf{Case 3: $G$ has no $C_3$ and no $C_4$.}
We now use the exact extremal tables introduced above.
Since
\[
|E(G)|=2n-4>\alpha_n=\operatorname{ex}(n;\{C_3,C_4,C_6\}),
\]
$G$ cannot be $\{C_3,C_4,C_6\}$-free. As we are already assuming $G$ has no $C_3$ and no $C_4$, it follows that $G$ contains a $6$-cycle. Hence
\[
6\in\mathcal C(G).
\]

We now subdivide Case~3.

\smallskip
\noindent\textbf{Case 3a: $G$ contains a $5$-cycle.}
Then $5,6\in\mathcal C(G)$, and so
\[
S(G)\ge \frac15+\frac16=\frac{11}{30}>\frac14.
\]

\smallskip
\noindent\textbf{Case 3b: $G$ has no $5$-cycle.}
Then $G$ has no $C_4$ and no $C_5$; moreover, from the previous paragraph it already has a $C_6$.
Now use the second exact extremal table:
\[
|E(G)|=2n-4>\beta_n=\operatorname{ex}(n;\{C_4,C_8\}).
\]
Therefore $G$ cannot be $\{C_4,C_8\}$-free. Since we are already assuming $G$ has no $C_4$, it follows that $G$ contains an $8$-cycle. Hence
\[
6,8\in\mathcal C(G),
\]
and therefore
\[
S(G)\ge \frac16+\frac18=\frac7{24}>\frac14.
\]

In all cases we have proved $S(G)\ge \frac14$.
Since equality is attained by $K_{2,n-2}$, the theorem follows.
\end{proof}

\begin{corollary}[New exact range on the $a=2$ complete-bipartite line]
For every $24\le n\le 39$,
\[
\min\{S(G): |V(G)|=n,\ |E(G)|=2n-4\}=\frac14,
\]
with minimiser $K_{2,n-2}$.
\end{corollary}

\begin{proof}
This is the $24\le n\le 39$ portion of Theorem~\ref{thm:round6-main}.
\end{proof}

\subsection*{6) ADVERSARIAL VERIFICATION}

I check the new argument against the main possible failure modes.

\begin{itemize}
\item \textbf{Dependence on Round~5.}
The only Round~5 inputs used are the already-settled range $7\le n\le 23$ and the formula $S(K_{2,n-2})=1/4$.
No earlier proof is reproved.

\item \textbf{Exactness of the extremal values.}
The ANU data page explicitly states that when the entry ``ne'' is not preceded by ``$\ge$'', the Tur\'an number is exact.
For every value of $\alpha_n$ and $\beta_n$ used here, the edge-number entry is exact.
For $\beta_n$ at $n=37,38,39$, the number of extremal graphs is recorded only as a lower bound, but the edge maximum itself is still exact, which is all the proof needs.

\item \textbf{Logical use of the forbidden-family bounds.}
In Case~3, the implication
\[
|E(G)|>\operatorname{ex}(n;\{C_3,C_4,C_6\}) \Longrightarrow C_6\subseteq G
\]
is valid because Case~3 has already excluded $C_3$ and $C_4$.
Likewise,
\[
|E(G)|>\operatorname{ex}(n;\{C_4,C_8\}) \Longrightarrow C_8\subseteq G
\]
is valid because Case~3b has already excluded $C_4$.

\item \textbf{Boundary of the new range.}
The inequalities are strict for every $24\le n\le 39$:
\[
2n-4>\alpha_n,
\qquad
2n-4>\beta_n.
\]
In particular, no equality case needs special handling.

\item \textbf{Why the method stops at $n=40$.}
The $C_6$ forcing still causes no difficulty near $n=40$, but the $C_8$ forcing argument uses exact values of $\operatorname{ex}(n;\{C_4,C_8\})$ only up to $n=39$ from the checked table.
Thus $n=40$ is the first value where the present exact-table proof no longer closes the argument.
This is a genuine limitation of the current method, not a logical gap in the proved range.
\end{itemize}

\subsection*{7) FINAL (EXACTLY ONE)}

\textbf{UNRESOLVED (BUT STRICTLY ADVANCED).}

Round~6 proves a strictly stronger theorem than Round~5:
for every $7\le n\le 39$, among all graphs with $n$ vertices and $2n-4$ edges,
\[
\min S(G)=\frac14,
\]
and the complete bipartite graph $K_{2,n-2}$ is a minimiser.
Thus the entire $a=2$ line is now settled through $n=39$.
The first unresolved value by the present method is $n=40$.
The full Erd\H{o}s--Hajnal minimiser question for arbitrary $(n,m)$ remains open.

\subsection*{8) COMPLETION ESTIMATE (MANDATORY)}

\textbf{COMPLETION: 94\%}

\subsection*{9) REFERENCES}

\begin{thebibliography}{9}

\bibitem{AfzalyMcKayData}
N.~Afzaly and B.~D.~McKay,
\emph{Extremal Graphs and Turan numbers},
Australian National University combinatorial data page.
Used here for the exact values of
$\operatorname{ex}(n;\{C_3,C_4,C_6\})$
and
$\operatorname{ex}(n;\{C_4,C_8\})$,
as well as for the page's convention explaining when an edge-number entry is exact.

\end{thebibliography}
